%% Elevark fra mattelab.no – laget automatisk av bygg.py. Rediger gjerne til eget bruk.
%% Kompiler med lualatex.
\documentclass[11pt,a4paper]{article}
\usepackage[norsk]{babel}
\usepackage[a4paper,margin=2.2cm,top=2cm,bottom=2.4cm]{geometry}
\usepackage{fontspec}
\usepackage{unicode-math}
\setmainfont{TeX Gyre Pagella}
\setsansfont{TeX Gyre Heros}[Scale=MatchLowercase]
\setmathfont{TeX Gyre Pagella Math}
\usepackage{xcolor}
\definecolor{blekk}{HTML}{16213A}
\definecolor{dempet}{HTML}{566077}
\definecolor{rute}{HTML}{CFC9B8}
\definecolor{m1}{HTML}{2F4AB0}
\definecolor{m2}{HTML}{99266D}
\usepackage[most]{tcolorbox}
\usepackage{qrcode}
\usepackage{needspace}
\usepackage{fancyhdr}
\usepackage[hidelinks]{hyperref}
\pdfvariable suppressoptionalinfo \numexpr 1+2+4+8+16+32+64+128+256+512\relax
\setlength{\parindent}{0pt}
\setlength{\parskip}{0.4em}
\pagestyle{fancy}
\fancyhf{}
\renewcommand{\headrulewidth}{0pt}
\fancyfoot[L]{\sffamily\footnotesize\color{dempet}mattelab.no · Faktorisering med produkt og sum · elevark uten fasit}
\fancyfoot[R]{\sffamily\footnotesize\color{dempet}Side \thepage}

\newcommand{\fase}[2]{\tcbox[on line, colframe=#2, colback=white, coltext=#2, boxrule=0.6pt, arc=2pt,
  boxsep=0pt, left=3pt, right=3pt, top=1.5pt, bottom=1.5pt, fontupper=\sffamily\bfseries\footnotesize]{#1}}
\newcommand{\steg}[4]{%
  \par\noindent{\color{rute}\rule{\linewidth}{0.4pt}}\par\nopagebreak
  \noindent\begin{minipage}[t]{2.1cm}\vspace{0pt}\fase{#1}{#2}\end{minipage}%
  \begin{minipage}[t]{\dimexpr\linewidth-2.1cm\relax}\vspace{0pt}#3\end{minipage}\par
  \vspace{#4}}
\newcommand{\del}[4]{%
  \par\noindent{\color{rute}\rule{\linewidth}{0.4pt}}\par\nopagebreak
  \noindent\begin{minipage}[t]{1.1cm}\vspace{0pt}\fase{#1}{#2}\end{minipage}%
  \begin{minipage}[t]{\dimexpr\linewidth-1.1cm\relax}\vspace{0pt}#3\end{minipage}\par
  \vspace{#4}}
\newcommand{\oppgave}[2]{\par\Needspace{12\baselineskip}\vspace{0.6em}%
  {\fontsize{14}{17}\selectfont\bfseries{\color{m1}#1}\enspace #2}\par\vspace{0.2em}}

\begin{document}
\begin{minipage}[t]{0.72\linewidth}\vspace{0pt}
{\sffamily\small\color{dempet}Elevark · 1T · mattelab.no}\par\vspace{0.3em}
{\fontsize{24}{28}\selectfont\bfseries Faktorisering med produkt og sum}\par\vspace{1em}
{\sffamily\small Navn: \rule{4.6cm}{0.4pt}\quad Klasse: \rule{1.5cm}{0.4pt}\quad Dato: \rule{1.9cm}{0.4pt}}
\end{minipage}\hfill
\begin{minipage}[t]{0.22\linewidth}\vspace{0pt}\raggedleft
\qrcode[height=2.3cm]{https://mattelab.no/utforskere/faktorisering.html}\par\vspace{0.2em}
{\sffamily\scriptsize\color{dempet}Utforskeren}
\end{minipage}

\vspace{1em}
{\sffamily\small\color{dempet}Utforskeren ligger på \url{https://mattelab.no/utforskere/faktorisering.html}}\par
Hver oppgave har tre steg. \textbf{Gjett} først, uten å bruke utforskeren. \textbf{Test} så i utforskeren. \textbf{Forklar} til slutt med egne ord hvorfor det blir slik.

\oppgave{1}{To tall med produkt og sum}
\steg{Gjett}{m1}{Finn to tall med produkt $6$ og sum $5$. Skriv $x^2 + 5x + 6$ som et produkt.}{2.2cm}
\steg{Test}{m2}{Vis $x^2 + 5x + 6$. Hvor mange tallpar har produkt $6$? Hvilket har sum $5$?}{2.6cm}
\steg{Forklar}{blekk}{Gang ut svaret ditt for å kontrollere. Forklar med rektangelet hvorfor produktet og summen må stemme.}{3.4cm}

\oppgave{2}{Negativt konstantledd}
\steg{Gjett}{m1}{Faktoriser $x^2 + x - 6$. Hvilke fortegn må de to tallene ha når produktet er negativt?}{2.2cm}
\steg{Test}{m2}{Vis $x^2 + x - 6$. Se på fortegnene i tabellen.}{2.6cm}
\steg{Forklar}{blekk}{Forklar hvorfor de to tallene må ha ulikt fortegn. Hvilke nullpunkter har uttrykket?}{3.4cm}

\oppgave{3}{Negativ sum, positivt produkt}
\steg{Gjett}{m1}{Faktoriser $x^2 - 7x + 12$.}{2.2cm}
\steg{Test}{m2}{Vis $x^2 - 7x + 12$. Finn alle tallparene med produkt $12$.}{2.6cm}
\steg{Forklar}{blekk}{Forklar hvorfor begge tallene må være negative. Hvordan ser du det på fortegnene til $b$ og $c$?}{3.4cm}

\oppgave{4}{Når ingen tall passer}
\steg{Gjett}{m1}{Prøv å faktorisere $x^2 + 4x + 5$. Går det?}{2.2cm}
\steg{Test}{m2}{Vis $x^2 + 4x + 5$. Hva står i tabellen? Hva er $b^2 - 4c$?}{2.6cm}
\steg{Forklar}{blekk}{Skriv uttrykket som $(x + 2)^2 + 1$. Bruk det til å forklare hvorfor uttrykket aldri blir $0$, og derfor ikke kan faktoriseres.}{3.4cm}

\oppgave{5}{Kvadratsetningene baklengs}
\steg{Gjett}{m1}{Faktoriser $x^2 + 6x + 9$ og $x^2 - 9$. Hva er spesielt med tallparene?}{2.2cm}
\steg{Test}{m2}{Vis $x^2 + 6x + 9$ og $x^2 - 9$. Se på formen til rektangelet.}{2.6cm}
\steg{Forklar}{blekk}{Hvilke kvadratsetninger kjenner du igjen? Forklar hvordan du kan se det på uttrykket uten å lete i en tabell.}{3.4cm}

\end{document}
